%% ma: L10-B14-0021
%% lop: 10
%% chuong: 5
%% bai: 14
%% dang: 10.14.2
%% loai: TN4
%% muc: TH
%% dap_an: "B"
%% nguon: "chuyen-de-cac-so-dac-trung-cua-mau-so-lieu-khong-ghep-nhom-toan-10-chuong-trinh-moi.pdf (GV. Phan Nhật Linh) (P04), Bài 03 Dạng 1 Phần I Câu 11"
%% kien_thuc: "Lớp 10 Bài 13 (tứ phân vị)"
%% trang_thai: chinh-thuc
%% ghi_chu: "Câu 14 gốc trùng hoàn toàn câu này, đã bỏ. Tự duyệt theo yêu cầu giáo viên 22/10/2026."
\begin{ex}
Điểm kiểm tra $9$ môn của một học sinh như sau:
\begin{center}
$6{,}5$\quad $8$\quad $9$\quad $8{,}5$\quad $7$\quad $7{,}5$\quad $6$\quad $7{,}25$\quad $9{,}5$
\end{center}
Các tứ phân vị $Q_1$, $Q_3$ của mẫu số liệu lần lượt là
\choice{$Q_1=6{,}75$; $Q_3=7{,}5$}{\True $Q_1=6{,}75$; $Q_3=8{,}75$}{$Q_1=7{,}5$; $Q_3=8{,}75$}{$Q_1=6{,}75$; $Q_3=8{,}5$}
\loigiai{Sắp xếp: $6$; $6{,}5$; $7$; $7{,}25$; $7{,}5$; $8$; $8{,}5$; $9$; $9{,}5$; $Q_2=7{,}5$.

Nửa dưới $6;6{,}5;7;7{,}25$ nên $Q_1=6{,}75$; nửa trên $8;8{,}5;9;9{,}5$ nên $Q_3=8{,}75$. Chọn B.}
\end{ex}
